% Folder 149, batch 04 — archivists' pages 61–80 (author's pp. 31–40, where visible).
% Pass: Opus 5.5 (claude-opus-5-5), 2026-10-03 — first pass, unchecked against the pages by a human.
%
% Notation fixed for this batch: \mathcal{E}_X for his curly E (the
% arithmetic fundamental group / extension attached to X); \mathcal{U} for
% his script U (the curve, the relative scheme over S); \overline{K} for the
% algebraic closure; \mathbb{Z}_\ell, T_\ell, \widehat{\mathbb{Z}}(1) as he
% writes them. The hand is fast pencil throughout; connective prose is often
% lost and is left \ill{}.
\documentclass[11pt,a4paper]{article}
\input{../preamble/grothendieck}

\folder{149}
\batch{4}
\pages{61}{80}
\foldertitle{[Autour de La "Longue Marche" à travers la théorie de Galois, pages 1 à 67] : notes manuscrites (s.d.).}
\dating{s.d. — le groupe « Autour de La "Longue Marche" à travers la théorie de Galois » (141 à 149) est daté [à partir de 1978]-1983}
\watermark{Édition de démonstration}

\begin{document}

\page{61}
\note{la phrase commence sur la page précédente, hors de ce lot.}

\ill{} tous nos beaux énoncés devraient être valables, non seulement pour un corps de base $K$ \uncertain{doté} de type fini sur $\mathbb{Q}$, mais dès que $K$ est ext.\ de type fini d'un corps cyclotomique (pas néc.\ fini sur $\mathbb{Q}$).

Nous pourrions définir les \emph{courbes de Poincaré} sur un corps alg.\ clos $\overline{K}$ de car.\ $0$ comme étant les courbes isomorphes à des rev.\ universels de courbes alg.\ anabéliennes sur $\overline{K}$ (deux courbes anabéliennes $\overline{\mathcal{U}}, \overline{\mathcal{V}}$ sur $\overline{K}$ définissant des rev.\ de Poincaré isomorphes, ssi $\exists$ un revêtement fini étale de l'une, isomorphe à un rev.\ fini étale de l'autre…).
\uncertain{Étant} donné une courbe de Poincaré $\widetilde{\overline{\mathcal{U}}}$ sur $\overline{K}$, \struck{telle} on définit canoniquement sa complétion $\widehat{\mathcal{U}}$, de telle façon que \uncertain{les automorphismes} de $(\widetilde{\overline{\mathcal{U}}}, \overline{K})$ (\uncertain{premiers}, \uncertain{induisant} l'identité sur $\overline{K}$) \uncertain{induisent} une \uncertain{action} de $\widehat{\mathcal{U}}$.
Ceci posé :

\page{62}
\note{p.~31 de l'auteur.}

\textbf{Conjecture B$'$.} Soient $K$ un corps de type fini sur $\mathbb{Q}$ (ou sur un corps cyclotomique, cela simplifie peut-être), $\overline{K}$ une clôture algébrique de $K$, $\mathcal{U}$ une courbe de Poincaré sur $\overline{K}$, de complétion $\widehat{\mathcal{U}} = \mathcal{X}$.
Considérons une action de $\Gamma = \Gamma_{\overline{K}/K}$ sur $\mathcal{U}$, compatible avec son action sur $\overline{K}$, d'où une action de $\Gamma$ sur $\mathcal{X}$. Ceci posé : il existe un point fixe et un seul de $\Gamma^{0}$ opérant sur $\mathcal{X}$

($\Gamma^{0}$, noyau du caractère cyclotomique $\Gamma \xrightarrow{\chi} \widehat{\mathbb{Z}}^{*}$).

La différence avec la conjecture B, pour celle-ci sous forme \uncertain{relative}, c'est qu'ici on ne suppose \emph{pas} que l'action de $\Gamma$ sur $\mathcal{U}$ normalise l'action d'un \emph{groupe} profini $\pi$, opérant \uncertain{librement} sur $\mathcal{U}$ de façon que $\mathcal{V} = \mathcal{U}/\pi$ soit une courbe alg.\ anabélienne sur $\overline{K}$.

\page{63}
\uncertain{Une circonstance} \ill{} les conjectures précédentes, quand on les applique : une situation où $K$ est engendré par un modèle $S$ de $K$ (disons, élémentaire anabélien), quand $\mathcal{U}_K$ provient d'une courbe relative $\mathcal{U}_S$ sur $S$ --- de sorte qu'on a un hom.\ de groupes
\[
(4)\qquad \mathcal{E}_{\mathcal{U}_S} \longrightarrow \mathcal{E}_S
\]
de noyau $\pi_{\mathcal{U}}$, dont $\mathcal{E}_{\mathcal{U}_K} \to \mathcal{E}_K$ se déduit par chgt de base i.e.\ par produit fibré

(5)

\begin{tikzcd}
\pi \arrow[d] \arrow[r, no head, "\sim"] & \pi \arrow[d] \\
\mathcal{E}_{\mathcal{U}_S} \arrow[d] & \mathcal{E}_{\mathcal{U}_K} \arrow[l] \arrow[d] \\
\mathcal{E}_S & \mathcal{E}_K \arrow[l, "\text{surj}"']
\end{tikzcd}

\note{au centre du carré, un petit signe \uncertain{$\varphi$} accompagné d'une flèche pointillée ; la flèche du bas porte « surj » en dessous.}

Ainsi, les sections de $\mathcal{E}_{\mathcal{U}_K}$ sur $\mathcal{E}_K$ correspondent aux \struck{certaines sections de} \add{relèvements continus} $\mathcal{E}_K \to \mathcal{E}_{\mathcal{U}_S}$ de l'hom.\ surjectif $\mathcal{E}_K \to \mathcal{E}_S$, et parmi ces relèvements, ceux qui sont \struck{\ill{}} \add{triviaux} sur le noyau de $\mathcal{E}_K \to \mathcal{E}_S$ correspondent exactement aux sections

\page{64}
\note{p.~32 de l'auteur.}

de $\mathcal{E}_{\mathcal{U}_S}$ sur $\mathcal{U}_S$. Nos conjectures impliquent donc qu'il y a exactement deux sortes de telles sections : 1°) celles qui correspondent à des pts de $\mathcal{U}_K/K$ i.e.\ à des sections rationnelles de $\mathcal{U}_S$ sur $S$ --- mais on va vérifier sans mal, sans doute, qu'une telle section rationnelle ne correspond effectivement à une section de \add{l'extension} (4), que si c'est une section \struck{rationnelle} \emph{régulière} (à vérifier tantôt). 2°) Celles correspondant à des $i \in I(\mathcal{U}_{\overline{K}})$ rationnels sur $K$, i.e.\ à une \emph{section} de $\widehat{\mathcal{U}}_S \smallsetminus \mathcal{U}_S = S'$ (étale fini sur $S$) sur $S$. Et il faudrait étudier encore : quelle condition une telle section définit un paquet non vide de scindages de (4) --- et comment déterminer exactement tous ces scindages.

Avant d'élucider ces deux points, un peu techniques, je voudrais voir

\page{65}
dans quelle mesure la conjecture A (ou B) faite dans le §, permet de reconstruire la catégorie des modèles élém.\ anabéliens sur $\mathbb{Q}$, \add{et} celle des extensions de type fini de $\mathbb{Q}$ et des modèles élém.\ anabéliens sur ceux-ci, en termes des \uncertain{groupes} \uncertain{extérieurs} \add{(ou topos \uncertain{galoisiens})} associés à partir bien sûr de la donnée fondamentale de $\Gamma_{\mathbb{Q}} = \Gamma_{\overline{\mathbb{Q}}/\mathbb{Q}_0}$, opérant extérieurement sur $\widehat{\pi}_{0,3}$, d'où déjà l'extension $\mathcal{E}_{0,3} = \mathcal{E}_{\mathcal{U}_{0,3}/\mathbb{Q}}$, où $\mathcal{U}_{0,3} = \mathbb{P}^1_{\mathbb{Q}} - \{0, 1, \infty\}$.
\note{l'indice de $\overline{\mathbb{Q}}/\mathbb{Q}_0$ est lu tel qu'écrit ; la lecture du second facteur est douteuse.}

Prenons les \add{\uncertain{donnée de} $\widehat{\pi}_{0,3}$-conj.\ de} sections de $\mathcal{E}_{0,3}$ sur $\Gamma_{\mathbb{Q}} = \Gamma$, \add{telles que la centralisateur du $\Gamma^{0}$ soit trivial (sections « \uncertain{admissibles} »)} \add{\uncertain{(\ill{} \ill{})}}, i.e.\ les « points » \struck{des} \add{topos} $B_{\widehat{\pi}_{0,3}, \Gamma_{\mathbb{Q}}}$ \ill{} sur $B_{\Gamma}$ --- on trouve un ensemble sur lequel $\Gamma$ opère (qui \ill{} \ill{} que $\mathcal{U}_{0,3}(\overline{\mathbb{Q}}_0)$, à \uncertain{un ens.}\ près). Pour \uncertain{tout} ens.\ fini $I$ de sections, \add{stable par $\Gamma$}, le foncteur

\page{66}
\note{p.~33 de l'auteur.}

« forage des trous » doit nous fournir un groupe extérieur $\pi_{0,3}(I)$, \add{du type $(0, 3 + \operatorname{card} I)$}, sur lequel $\Gamma$ opère (il vaut mieux peut-être utiliser le yoga introduit par Ihara des groupoïdes rigides --- donc on prend deux \uncertain{ou trois} bouts \uncertain{entrant} en trous $0, 1, \infty$ --- \ill{} donc l'équivalent groupoïdal de la \struck{sphère \ill{}} droite projective $\mathbb{P}^1_{\mathbb{Q}}$, on l'appelle la sphère \uncertain{arithmétique} --- qui correspond à un groupe \add{ext} à lacets \uncertain{non-fini} sur lequel $\Gamma$ opère --- en fait, \uncertain{ce n'est autre} que $\mathcal{E}_{K_1}$, où
\[
(6)\qquad K_1 = \mathbb{Q}(\text{\struck{$X$}}\, T_1)
\]
est l'extension transcendante pure \add{type} de degré $1$ de $\mathbb{Q}$.

Partant de (6), on construit de \uncertain{nouveau} l'équivalent groupoïdal de \struck{\ill{}} $\mathcal{U}_{0,3}$ \struck{(\ill{}} et on \uncertain{retrouve} comme

\page{67}
précédemment, \uncertain{pour savoir}, \uncertain{tant} des courbes de type $(0, \nu_2)$ sur $K_1$ (ou sur une extension finie de $K_1$) \uncertain{tant} des courbes relatives de type $(0, \nu_2)$ sur une courbe sur $\mathbb{Q}$ \struck{de type} (ou une extension finie de $\mathbb{Q}$, ou un rev.\ \struck{étale} étale fini d'une telle $\mathcal{U}_{0,\nu_1}$).
On procède de proche en proche pour construire finalement tous les $\mathcal{E}_K$ \uncertain{sur} $\mathcal{E}_{\mathbb{Q}}$ (car tt corps \struck{ext.}\ de type fini de $\mathbb{Q}$, est ext.\ finie d'une ext.\ transcendante pure) et tous les modèles élémentaires, où à chaque cran la fibration est soit une courbe de genre $0$, soit un rev.\ étale fini d'une telle fibration particulière. Sauf erreur, \struck{ça} fait assez pour avoir un syst.\ fondamental de voisinages de tt point d'une $X$ lisse sur un $K$

\page{68}
\note{p.~34 de l'auteur.}

et de reconstituer en principe les schémas lisses \uncertain{sur} des $K$, par recollements de tels morceaux \uncertain{avec} des « \uncertain{immersions} ouvertes ». Mais je suis certain que pour faire une telle description, il me faudra développer un \uncertain{langage} géométrique qui colle mieux à l'intuition géométrique, que les \uncertain{sempiternelles} extensions de groupes profinis… ou actions extérieures, et où les points rigides (\uncertain{ou} \uncertain{cuspidaux} dans des clôtures algébriques de corps) jouent un rôle prépondérant. Il me faudra y revenir dessus --- et en même temps, expliciter les topos étales (pas seulement le « morceau $K(\pi,1)$ ») \uncertain{tout} entiers des schémas décrits ici par des extensions.

\page{69}
Reprenons le cas de $\mathcal{U} = \mathcal{U}_S$ schéma relatif sur $S$, « élémentaire » sur $S$ --- à fibres successives anabéliennes (s'il le faut) ou du moins à $\pi_1$ non nul, $S$ étant lui-même (pour fixer les idées) lisse sur $\mathbb{Q}$, irréductible, corps de fonctions $K$, et reprenons le diagramme 5. Considérons une section rationnelle \add{$f$} de $\mathcal{U}$ sur $S$, définissant une section de $\mathcal{E}_{\mathcal{U}_K}$ sur $\mathcal{E}_K$ --- ou, ce qui revient au même, un relèvement \add{\uncertain{de section}} de l'hom.\ surjectif $\mathcal{E}_K \to \mathcal{E}_S$ en $\mathcal{E}_K \to \mathcal{E}_{\mathcal{U}_S}$ \add{(composé de $\mathcal{E}_K \to \mathcal{E}_{\mathcal{U}_K}$ avec $\mathcal{E}_{\mathcal{U}_K} \to \mathcal{E}_{\mathcal{U}_S}$)}. Je veux montrer que $f$ est partout définie, i.e.\ une \struck{\uncertain{section partout définie}} \uncertain{section} de $\mathcal{U}_S$ sur $S$, ssi la section de $\mathcal{E}_{\mathcal{U}_K}$ sur $\mathcal{E}_K$ provient d'une section de $\mathcal{E}_{\mathcal{U}_S}$ sur $\mathcal{E}_S$, i.e.\ ssi le relèvement $\mathcal{E}_K \to \mathcal{E}_{\mathcal{U}_S}$ s'annule sur

\page{70}
\note{p.~35 de l'auteur.}

le noyau de $\mathcal{E}_K \to \mathcal{E}_S$.

Notons que cette dernière condition est une condition « de codimension 1 sur $S$ » --- de façon plus précise, si $Z$ est un sous-schéma fermé de $S$ de codim $\geqslant 2$, \struck{alors} alors, posant $S' = S \smallsetminus Z$, on a $\pi_1(S') \xrightarrow{\sim} \pi_1(S) = \mathcal{E}_S$ par le « théorème de pureté » --- donc le noyau de \struck{$\pi_1$} $\mathcal{E}_K \to \mathcal{E}_S$ est le même que celui de $\mathcal{E}_K \to \mathcal{E}_{S'}$, or, \uncertain{si on prend} (comme $S'$ \uncertain{n'est} pas un « modèle ») que le sous-groupe \struck{fermé} engendré par les noyaux des $\mathcal{E}_K \to \mathcal{E}_{S'_i}$, où les $S'_i$ sont des ouverts « modèles » qui recouvrent $S'$. \uncertain{Si} donc la condition envisagée est remplie relativement aux $S'_i$ (qui pourraient \struck{\ill{}} ne pas recouvrir $S$, seulement $S'$) --- ce qui est aussi signifie que la section

\page{71}
rationnelle envisagée est \uncertain{morphique} sur les $S'_i$, i.e.\ sur $S'$ --- alors elle est vérifiée \struck{relativement} à $S$ --- ce qui est aussi \emph{signifier} que la section est \uncertain{morphique} sur $S$. Donc, pour se rassurer, il faudrait s'assurer a priori qu'une section de $\mathcal{U}_S$, sur $S'$ se prolonge automatiquement \struck{en $\mathcal{U}$} en une section de $\mathcal{U}_S$ sur $S$. \struck{De} De proche en proche, on est ramené au cas d'une courbe relative $\mathcal{U}_S = X_S - T$, $X$ lisse \add{sur $S$} \struck{de} dim.\ relative $1$, $T$ fini étale sur $S$, \struck{\ill{}} \uncertain{ou} $T$ décomposé sur $S$. Si $X$ de genre relatif $\geqslant 1$, on sait (th.\ de Weil) que la section se prolonge en une section de $X$, soit $D$ l'image inverse de $T$, c'est un diviseur sur $S$, dont la trace sur $S' = S \smallsetminus Z$ est nul, donc (comme $\operatorname{codim}(Z, S) \geqslant 2$) il est nul, OK.

\page{72}
\note{p.~36 de l'auteur.}

(9)

\begin{tikzcd}[column sep=small]
\mathcal{U}_S \arrow[d, no head] & \mathcal{U}_{\mathcal{O}_x} \arrow[l, no head] \arrow[d, no head] & \mathcal{U}_K \arrow[l, no head] \arrow[d, no head] & \mathcal{U}_{\overline{K}} \arrow[l, no head] \arrow[d, no head] & \widetilde{\mathcal{U}}_S \arrow[l, no head] \\
S & \mathcal{O}_x \arrow[l, no head] & K \arrow[l, no head] & \overline{K} \arrow[l, no head] &
\end{tikzcd}

(10)

\begin{tikzcd}
\mathcal{E}_{\mathcal{U}_S} \arrow[d] & \mathcal{E}_{\mathcal{U}_{\mathcal{O}_x}} \arrow[l] \arrow[d] & \mathcal{E}_{\mathcal{U}_K} \arrow[l] \arrow[d] \\
\mathcal{E}_S & \mathcal{E}_{\mathcal{O}_x} \arrow[l] & \mathcal{E}_K \arrow[l] \arrow[ul, dashed, "\varphi_x"']
\end{tikzcd}

\note{dans (9), les traits sont sans flèches ; $\mathcal{O}_x$ est souligné sur la page (anneau local en $x$).}

\uncertain{avec} des carrés cartésiens, et des flèches horizontales surjectives. L'hom.\ $\mathcal{E}_K \leftarrow \mathcal{E}_{\mathcal{U}_S}$ \note{sic : sens de la flèche tel qu'écrit.} est composé d'un relèvement $\mathcal{E}_K \to \mathcal{E}_{\mathcal{U}_{\mathcal{O}_x}}$ de $\mathcal{E}_K \to \mathcal{E}_{\mathcal{O}_x}$, \uncertain{avec} l'hom.\ \struck{dernier} canonique $\mathcal{E}_{\mathcal{U}_{\mathcal{O}_x}} \to \mathcal{E}_{\mathcal{U}_S}$. (\struck{Les} Les relèvements $\mathcal{E}_K \to \mathcal{E}_{\mathcal{U}_{\mathcal{O}_x}}$ correspondent \uncertain{biunivoquement} \uncertain{aux} sections de $\mathcal{E}_{\mathcal{U}_K}$ sur $\mathcal{E}_K$, \uncertain{comme aux} relèvements de $\mathcal{E}_K \to \mathcal{E}_S$ en $\mathcal{E}_K \to \mathcal{E}_{\mathcal{U}_S}$…)

\marginal{NB La quantité de $\mathcal{E}_K$ par \uncertain{le} \ill{} engendré par $I_x$ \uncertain{n'est autre} que $\mathcal{E}_{\mathcal{O}_x} = \pi_1(\operatorname{Spec} \mathcal{O}_x)$.}
\note{note marginale écrite en oblique dans la marge gauche, en bas de page ; lecture en grande partie douteuse.}

Ceci dit, j'ai envie de prouver que $\varphi_x : \mathcal{E}_K \to \mathcal{E}_{\mathcal{U}_{\mathcal{O}_x}}$ s'annule \struck{sur} \uncertain{sur} $I_x$, i.e.\ provient d'une section de $\mathcal{E}_{\mathcal{U}_{\mathcal{O}_x}}$ sur $\mathcal{O}_x$ ssi la section rationnelle correspondante de $\mathcal{U}_S/S$ est définie en $x$. Ceci impliquera

\page{73}
l'assertion précédente (que \struck{la} la section $\varphi$ de $\mathcal{E}_{\mathcal{U}_K}$ sur $\mathcal{E}_K$ provient d'une section de $\mathcal{E}_{\mathcal{U}_S}$ sur $\mathcal{E}_S$, ssi la section rationnelle correspondante est \uncertain{morphique}).

Mais il s'\emph{agit} ici d'un énoncé en fait purement géométrique, que j'ai envie de reformuler sous forme plus générale :

\textbf{Théorème} \struck{\ill{}}\uncertain{1)}. Soit $T$ un trait (spectre d'un anneau \uncertain{de valuation discrète}), $\mathcal{U}$ un schéma relatif « élémentaire » sur $T$, anabélien$^{*}$, $K$ le corps des fonctions de $T$. On choisit un rev.\ universel $\widetilde{\mathcal{U}}$ de $\mathcal{U}$, d'où une clôture alg.\ $\overline{K}$ de $K$, et on considère l'extension $\mathcal{E}_{\mathcal{U}} = \pi_1(\mathcal{U}; \widetilde{\mathcal{U}})$ de $\mathcal{E}_K = \pi_1(K, \overline{K}) \simeq \operatorname{Gal}(\overline{K}/K)$ par $\pi = \pi_1(\mathcal{U}_{\overline{K}}, \widetilde{\mathcal{U}})$. \struck{Considérons \add{On a donc} une section pt $K$-rationnel de $\mathcal{U}_K$} une \add{\uncertain{suite}} extension de groupes profinis

\marginal{$^{*}$ anabéliennement \uncertain{inutile} --- il suffit que les fibres des fibrations élémentaires de $\mathcal{U}$ ne soient pas du type $(0,0)$ ou $(0,1)$, \uncertain{de} \uncertain{dimension} relative $1$ \uncertain{et} \uncertain{à} $\pi_1$ nul.}
\note{note marginale en oblique dans la marge gauche, renvoyée par l'astérisque après « anabélien » ; lecture en partie douteuse.}

\struck{d'où}

\begin{tikzcd}
\mathcal{E}_{\mathcal{U}} \arrow[d, "\text{surj}"'] & \mathcal{E}_{\mathcal{U}_K} \arrow[l, "\text{surj}"'] \arrow[d, "\text{surj}"] \\
\mathcal{E}_S & \mathcal{E}_K = \operatorname{Gal}(\overline{K}/K) \arrow[l, "\text{surj}"]
\end{tikzcd}

où $\mathcal{E}_S$ s'identifie au quotient de $\mathcal{E}_K$ par le sous-groupe \uncertain{normal} fermé engendré par un groupe

\note{dans le diagramme, $\mathcal{E}_S$ (et non $\mathcal{E}_T$) est tel qu'écrit.}

\page{74}
\note{p.~37 de l'auteur.}

d'inertie $I_{x} \simeq T_{\infty}(\overline{K})$, cf.\ plus haut. Soit $f_K$ \uncertain{$\in$} $K$ un pt de $\mathcal{U}_K$ rat./$K$, d'où une section $\psi = \psi_{f_K}$ de $\mathcal{E}_K$ sur $\mathcal{E}_{\mathcal{U}_K}$. Ceci posé, les conditions suivantes sont équivalentes
\begin{itemize}
\item[a)] $f_K$ se prolonge en une section de $\mathcal{U}$ sur $S$
\item[b)] $\psi$ provient d'une section de $\mathcal{E}_{\mathcal{U}}$ sur $\mathcal{E}_S$
\item[b$'$)] le composé $\mathcal{E}_K \xrightarrow{\psi} \mathcal{E}_{\mathcal{U}_K} \to \mathcal{E}_{\mathcal{U}}$ s'annule sur $I_x$
\end{itemize}
\struck{\ill{}}

\note{« section $\psi$ de $\mathcal{E}_K$ sur $\mathcal{E}_{\mathcal{U}_K}$ » : ordre des termes tel qu'écrit ; $S$ et $T$ désignent ici le même trait.}

L'équivalence de b) et b$'$), et qu'elles soient impliquées par a), est claire. C'est l'implication b) $\Longrightarrow$ a) qui demande une démonstration. On est ramené \uncertain{aisément}, par descente, au cas \struck{\ill{}} où $T$ est strictement local (donc $\mathcal{E}_K = \operatorname{Gal}(\overline{K}/K)$ est réduit à son ss-groupe d'inertie\add{, et $\mathcal{E}_S = (1)$}). On est ramené de proche en proche au cas où $\mathcal{U}/S$ est une courbe relative élémentaire, \struck{de genre $g \geqslant 1$,} $\mathcal{U} = X \smallsetminus T$, $X$ lisse \add{et propre}. Alors $f$ se prolonge en une section $f$ \add{de} \struck{rationnelle} $X$ sur $S$, et la conclusion \uncertain{signifie} que $f(S) \subset \mathcal{U}$. Donc on est ramené finalement \uncertain{à} --

\textbf{Lemme.} Soit $X$ schéma projectif lisse de dim.\ relative $1$ connexe sur $S$ trait strictement local,

\page{75}
Soit $T \subset X$ sous schéma fini étale sur $S$, donc $T \simeq I_S$, $I$ ensemble fini, et soit $\mathcal{U} = X \smallsetminus T$ (donc $T$ est défini par une famille $(g_i)_{i \in I}$ de sections disjointes de $X$ sur $S$)\add{. Si $g$ est le genre relatif, $\nu = \operatorname{card} I$, on suppose $(g, \nu) \neq (0,1)$}. Soit \uncertain{fixé} $i_0 \in I$, $f$ une section \add{de $X/S$} \struck{disjointe} distincte des $g_i$, et telle que \struck{$f \mid k(s) = g_{i_0}$} $f$ et $g_{i_0}$ coïncident en $s$ (pt fermé de $S$). \struck{Considérons} \add{Donc} \add{si $\mathcal{Y} = S \smallsetminus \{s\}$,} \struck{\ill{} \ill{} morphisme} on a donc un morphisme $\mathcal{Y} \to \mathcal{U}$, d'où
\[
\pi_1(\mathcal{Y}) \longrightarrow \pi_1(\mathcal{U}).
\]
Je dis que cet hom.\ \add{si $\nu \geqslant 2$}, n'est pas trivial, \struck{\ill{}} et même, que pour \add{$\ell$ donné}
\[
T_\ell(\widehat{K}^{*}) \simeq H_1(\mathcal{Y}, \mathbb{Z}_\ell) \longrightarrow H_1(\mathcal{U}, \mathbb{Z}_\ell) \;\bigl[\simeq H_1(\mathcal{U}_s, \mathbb{Z}_\ell)\bigr]
\]
n'est pas trivial (pour \struck{\ill{}} $\ell$ nombre premier distinct de la caractéristique résiduelle).

\note{« $T_\ell(\widehat{K}^{*})$ » est écrit dans la marge gauche, devant $H_1(\mathcal{Y}, \mathbb{Z}_\ell)$ ; $[H_1(\mathcal{U}_s, \mathbb{Z}_\ell)]$ est écrit sous $H_1(\mathcal{U}, \mathbb{Z}_\ell)$, relié par une flèche verticale.}

Comme la section rationnelle de $J^{1}_{X/S}$ définie par $f$ est régulière, on voit que le composé de l'hom.\ envisagé avec $H_1(\mathcal{U}, \mathbb{Z}_\ell) \to H_1(\widehat{J}^{1}_{X/S}, \mathbb{Z}_\ell)$ est nul --- i.e.\ le $H_1(\mathcal{Y}, \mathbb{Z}_\ell)$ s'envoie dans la partie torique de $H_1(\mathcal{U}, \mathbb{Z}_\ell)$ ($\xleftarrow{\sim} H_1(\mathcal{U}_s, \mathbb{Z}_\ell)$), qui est canoniquement isomorphe à $\mathbb{T}_\ell^{I} / \mathbb{T}_\ell$.
[NB cette partie est nulle si $\operatorname{card} I = 1$, et

\note{la lettre notée ici $J^{1}_{X/S}$ (Jacobienne, ou torseur de degré~$1$) est une lecture douteuse d'un $J$ cursif ; de même $\mathbb{T}_\ell$, écrit en gras, pour le module de Tate de $\mathbb{G}_m$.}

\page{76}
\note{p.~38 de l'auteur.}

dans ce cas le critère homologique serait insuffisant…) Il faudrait donc calculer cet \struck{élément de} \add{homomorphisme}
\[
(11)\qquad T_\ell\,(\simeq H_1(\mathcal{Y}, \mathbb{Z}_\ell)) \longrightarrow \mathbb{T}_\ell^{I} / T_\ell
\]
pour constater qu'il n'est pas nul dans le cas envisagé, \add{$\nu \geqslant 2$} (et traiter séparément le cas $\nu = 1$).

\note{le passage qui suit, jusqu'à « isogénies », est barré de trois grands traits obliques sur la page.}

\struck{[On va faire des calculs intrinsèques, \uncertain{avec}}
\[
\text{\struck{$\operatorname{Ker}\bigl(J^{1}_{\mathcal{U}/S} \to J^{1}_{X/S}\bigr) \simeq \mathbb{G}_m^{I} / \mathbb{G}_m$}}
\]
\[
\text{\struck{$H_1(\mathcal{U}, \mathbb{Z}_\ell) \xrightarrow{\sim} H_1(J^{1}_{\mathcal{U}/S}, \mathbb{Z}_\ell) \simeq \varprojlim {}_{\ell^n}(J^{0}_{\mathcal{U}/S})$}}
\]
\struck{(\uncertain{isom.\ de Kummer}). On considère donc l'image inverse par $\mathcal{Y} \to J^{1}_{\mathcal{U}/S}$ (composé de $\mathcal{Y} \to \mathcal{U} \to J^{1}_{\mathcal{U}/S}$) des isogénies}

Je vais deviner le résultat : Soit \struck{\ill{}} $x = g_{i_0}(s)$, $A = \mathcal{O}_{X,x}$, $V$ l'anneau de $S$, $J_{i_0}$ l'idéal de l'hom.\ $A \xrightarrow{g_{i_0}^{*}} V$ associé à $g_{i_0}$, c'est donc un idéal inversible de $A$ --- soit de même $J_f$ l'idéal associé à $f^{*} : A \to V$, et considérons $g_{i_0}^{*}(J_f)$, c'est un idéal de $V$ engendré par un générateur, et comme $g_{i_0} \neq f$, on voit que cet idéal n'est pas nul. Soit $n = \operatorname{long} V / g_{i_0}^{*}(J_f)$ \struck{$=$}, est entier

\page{77}
\uncertain{et} \uncertain{dépend} pas \uncertain{que} \uncertain{intrinsèquement} \uncertain{des} \uncertain{valeurs} de $g_{i_0}$ et $f$, car
\[
V / g_{i_0}^{*}(J_f) \simeq A / (J_{g_{i_0}} + J_f)
\]
--- \uncertain{donc} $n$ est la longueur du schéma artinien local, intersection des \add{\uncertain{sous-schémas}} $f(S)$ et $g_{i_0}(S)$ de $X$ --- on peut aussi l'interpréter comme une multiplicité d'intersection. Ceci posé, je présume que l'hom.
\[
T_\ell \longrightarrow T_\ell^{I} / T_\ell
\]
est le produit par $n$ de l'injection canonique $T_\ell \to T_\ell^{I}$, correspondant à l'indice $i_0$. Il faudrait que je demande à Carlos de me le confirmer.

Reste le cas $\nu = 1$, qui semble demander un traitement séparé. \struck{Il} \add{Il} suffit à vérifier (pour les groupes fondamentaux premiers à $p$) c'est que l'hom.\ ext.\ $\pi_1(\mathcal{Y}) \to \pi_1(\mathcal{U}) \simeq \pi_1(\mathcal{U}_s)$ est égal à \struck{\ill{}} $k_{i_0} \circ (n \cdot \mathrm{id}_T)$, où $k_{i_0}$ est

\marginal{Ceci dit, \uncertain{sans} \uncertain{indépendant} \uncertain{de} \uncertain{la} \uncertain{valeur} \uncertain{de} $\nu$ !}
\note{note marginale en oblique, à gauche du paragraphe sur le cas $\nu = 1$ ; lecture largement douteuse. « Carlos » : prénom tel qu'écrit.}

\page{78}
\note{p.~39 de l'auteur.}

l'hom.\ « local »
\[
(12)\qquad k_{i_0} : T \;(= T_{\infty}(k)) \longrightarrow \pi_1(\mathcal{U}_k)
\]
associé à l'indice $i_0$. Je vais admettre ce point, qui ne peut guère être difficile.

Pour terminer ce \uncertain{n°}, je veux encore étudier, dans la situation d'une $\mathcal{U}$ \struck{élémentaire} \add{courbe} relative sur un $S$ \uncertain{avec} \struck{\ill{}} $\mathcal{U} = X \smallsetminus T$, \add{$X$ lisse et propre sur $S$,} $T$ fini étale, une section $g_i$ donnée de $T$ sur $S$, les « sections de \uncertain{2ème espèce} » de l'extension
\[
1 \longrightarrow \pi \longrightarrow \mathcal{E}_{\mathcal{U}} \longrightarrow \mathcal{E}_S \longrightarrow 1
\]
associées à $i = i_0$ --- qui dépend \add{d'une classe de $\pi$-conj.\ de} un \add{(\uncertain{ss}-)}groupe \uncertain{de} \uncertain{lacets} Lie \uncertain{dans} $\pi$. (On suppose ici qu'on a bien une telle suite exacte, i.e.\ que $\pi_2(S) \to \pi_1(\text{fibre})$ est nul, ce qui sera le cas si $\pi_2(S) = 0$, p.\ ex.\ (je présume) si on est dans le cas d'un modèle élémentaire \uncertain{au-dessus} d'un corps de car.\ zéro (la restriction de car.\ nulle étant sans doute inutile, si on a bien

\marginal{$^{*}$ En fait, \struck{on n'a pas \uncertain{besoin} que le centre de $\pi$ soit trivial pour cette construction --- c'est un fait général pour \uncertain{toute} situation \uncertain{fibrée}, \uncertain{qui} \uncertain{fournit} une ext.\ \uncertain{des} $\pi_1$ de la base par le $\pi_1$ de la fibre, \uncertain{comme}} \uncertain{en} \uncertain{tous} \uncertain{cas}, \uncertain{il} \uncertain{suffit} \uncertain{que} $\pi_1(S) \to \pi$ \uncertain{triviale} \ill{} $\mathcal{E}_S = \pi_1(S)$ opère \uncertain{extérieurement} sur $\pi = \pi_1(\text{fibre géom.})$, d'où une extension \ill{} $\pi$ \uncertain{associée} à cette \uncertain{triviale}… \uncertain{celle} \ill{} $\pi_1(\mathcal{U})$ \ill{}… \uncertain{Il faudrait} \uncertain{regarder} \uncertain{ceci} \uncertain{de} \uncertain{plus près} !}
\note{longue note marginale en oblique dans la marge gauche, dont un premier bloc est barré de traits obliques ; lecture très incertaine, et la place de l'astérisque de renvoi dans le texte n'a pas été retrouvée.}

\page{79}
aux groupes fondamentaux premiers aux car.\ résiduelles…) Par un argument général, le normalisateur de $L_i$ dans $\mathcal{E}_{\mathcal{U}}$ s'envoie \emph{sur} $\mathcal{E}_S$, on trouve donc des scindages par cette extension, qu'on peut regarder comme une extension
\[
(13)\qquad 1 \longrightarrow T \longrightarrow N(L_i) \longrightarrow \mathcal{E}_S \longrightarrow 1 .
\]
La classe d'isomorphie est un élément \struck{\ill{}}
\[
(14)\qquad c_i \in H^2(\mathcal{E}_S, T), \qquad T = \varprojlim \mu_n ,
\]
que je me propose d'étudier. On a, si $S$ est un $K(\pi, 1)$
\[
(15)\qquad H^2(\mathcal{E}_S, T) \simeq H^2(S, T) \simeq \varprojlim_n H^2(S, \mu_n) ;
\]
d'ailleurs on a une suite exacte de Kummer (\uncertain{où} $\operatorname{Pic}(S) = H^1(S, \mathbb{G}_m)$)
\[
(16)\qquad 0 \to \operatorname{Pic}(S)/n \to H^2(S, \mu_n) \to {}_nH^2(S, \mathbb{G}_m) \to 0
\]
d'où par passage à la limite
\[
(17)\qquad 0 \to \operatorname{Pic}(S)^{\wedge} \to H^2(S, T) \to \varprojlim {}_nH^2(S, \mathbb{G}_m)
\]
\note{sous $\operatorname{Pic}(S)^{\wedge}$ : « $\overset{\text{déf}}{=} \varprojlim \operatorname{Pic}/n\operatorname{Pic}$ ».}

Dans le cas où $S$ est un schéma élémentaire sur un corps de \struck{car.\ \uncertain{0}} \ill{}

\note{dans (14), le « $T$ » est relié par un signe d'égalité vertical à $\varprojlim \mu_n$ écrit dessous.}

\page{80}
\note{p.~40 de l'auteur.}

de type fini sur $\mathbb{Q}$, $\operatorname{Pic}(S)$ est un $\mathbb{Z}$-\struck{groupe} \add{module} de type fini \add{(par Mordell--Weil--Néron)}, donc l'hom.
\[
(18)\qquad \operatorname{Pic}(S) \longrightarrow \operatorname{Pic}(S)^{\wedge} \hookrightarrow H^2(S, T)
\]
est \emph{injectif}.

Sans nous soucier de cette condition, \struck{pour} considérons le cas général --- je dis que la classe \struck{\ill{}} $c$ \add{(14)} est dans l'image de (18), de façon précise que c'est l'image de l'élément
\[
(19)\qquad \xi_i \in \operatorname{Pic}(S)
\]
classe du faisceau conormal (\uncertain{ou} normal, \uncertain{le} \uncertain{diable} \uncertain{sait}) de $X$ le long de $g_i$. Principe d'une vérification : on prend le complété formel de $X$ le long de $g_i(S)$ \struck{\ill{}} --- \uncertain{géométriquement} --- un \uncertain{voisinage} \uncertain{tubulaire} \uncertain{de} $g_i(S)$, et on interprète la suite exacte (13) comme la suite exacte d'homotopie de ce \struck{schéma} topos \uncertain{fibré}, au-dessus de \struck{\ill{}} $S$. On \uncertain{est} donc à préciser une histoire d'\uncertain{orbites}…

\note{la page s'arrête ici ; la vérification esquissée se poursuit vraisemblablement dans le lot suivant.}

\end{document}
